% Options for packages loaded elsewhere
\PassOptionsToPackage{unicode}{hyperref}
\PassOptionsToPackage{hyphens}{url}
\PassOptionsToPackage{dvipsnames,svgnames,x11names}{xcolor}
%
\documentclass[
  10.5pt,
  a4paper,
]{article}
\usepackage{amsmath,amssymb}
\usepackage{iftex}
\ifPDFTeX
  \usepackage[T1]{fontenc}
  \usepackage[utf8]{inputenc}
  \usepackage{textcomp} % provide euro and other symbols
\else % if luatex or xetex
  \usepackage{unicode-math} % this also loads fontspec
  \defaultfontfeatures{Scale=MatchLowercase}
  \defaultfontfeatures[\rmfamily]{Ligatures=TeX,Scale=1}
\fi
\usepackage{lmodern}
\ifPDFTeX\else
  % xetex/luatex font selection
  \setmainfont[]{DejaVu Serif}
  \setsansfont[]{DejaVu Sans}
  \setmonofont[]{DejaVu Sans Mono}
\fi
% Use upquote if available, for straight quotes in verbatim environments
\IfFileExists{upquote.sty}{\usepackage{upquote}}{}
\IfFileExists{microtype.sty}{% use microtype if available
  \usepackage[]{microtype}
  \UseMicrotypeSet[protrusion]{basicmath} % disable protrusion for tt fonts
}{}
\makeatletter
\@ifundefined{KOMAClassName}{% if non-KOMA class
  \IfFileExists{parskip.sty}{%
    \usepackage{parskip}
  }{% else
    \setlength{\parindent}{0pt}
    \setlength{\parskip}{6pt plus 2pt minus 1pt}}
}{% if KOMA class
  \KOMAoptions{parskip=half}}
\makeatother
\usepackage{xcolor}
\usepackage[margin=2.2cm]{geometry}
\usepackage{longtable,booktabs,array}
\usepackage{calc} % for calculating minipage widths
% Correct order of tables after \paragraph or \subparagraph
\usepackage{etoolbox}
\makeatletter
\patchcmd\longtable{\par}{\if@noskipsec\mbox{}\fi\par}{}{}
\makeatother
% Allow footnotes in longtable head/foot
\IfFileExists{footnotehyper.sty}{\usepackage{footnotehyper}}{\usepackage{footnote}}
\makesavenoteenv{longtable}
\usepackage{graphicx}
\makeatletter
\def\maxwidth{\ifdim\Gin@nat@width>\linewidth\linewidth\else\Gin@nat@width\fi}
\def\maxheight{\ifdim\Gin@nat@height>\textheight\textheight\else\Gin@nat@height\fi}
\makeatother
% Scale images if necessary, so that they will not overflow the page
% margins by default, and it is still possible to overwrite the defaults
% using explicit options in \includegraphics[width, height, ...]{}
\setkeys{Gin}{width=\maxwidth,height=\maxheight,keepaspectratio}
% Set default figure placement to htbp
\makeatletter
\def\fps@figure{htbp}
\makeatother
\setlength{\emergencystretch}{3em} % prevent overfull lines
\providecommand{\tightlist}{%
  \setlength{\itemsep}{0pt}\setlength{\parskip}{0pt}}
\setcounter{secnumdepth}{-\maxdimen} % remove section numbering

\usepackage{fancyhdr}
\pagestyle{fancy}
\fancyhf{}
\fancyhead[L]{\small Cooperative payload transport --- concept guide}
\fancyhead[R]{\small\thepage}
\renewcommand{\headrulewidth}{0.3pt}
\usepackage{xcolor}
\definecolor{acc}{HTML}{2A78D6}
\usepackage{sectsty}
\allsectionsfont{\color{acc}\sffamily}
\setlength{\parskip}{4pt}
\usepackage{float}
\let\origfigure\figure
\let\endorigfigure\endfigure
\renewenvironment{figure}[1][2] {\expandafter\origfigure\expandafter[H]} {\endorigfigure}
\ifLuaTeX
  \usepackage{selnolig}  % disable illegal ligatures
\fi
\IfFileExists{bookmark.sty}{\usepackage{bookmark}}{\usepackage{hyperref}}
\IfFileExists{xurl.sty}{\usepackage{xurl}}{} % add URL line breaks if available
\urlstyle{same}
\hypersetup{
  pdftitle={Cooperative Payload Transport by a Robot Formation},
  pdfauthor={Companion to the benchmark and the PA-VFH manuscript},
  colorlinks=true,
  linkcolor={acc},
  filecolor={Maroon},
  citecolor={Blue},
  urlcolor={acc},
  pdfcreator={LaTeX via pandoc}}

\title{Cooperative Payload Transport by a Robot Formation}
\usepackage{etoolbox}
\makeatletter
\providecommand{\subtitle}[1]{% add subtitle to \maketitle
  \apptocmd{\@title}{\par {\large #1 \par}}{}{}
}
\makeatother
\subtitle{Concept guide: learn every idea, read every graph}
\author{Companion to the benchmark and the PA-VFH manuscript}
\date{November 2026}

\begin{document}
\maketitle

{
\hypersetup{linkcolor=black}
\setcounter{tocdepth}{2}
\tableofcontents
}
This guide teaches every idea the project uses, from the ground up, and
then shows how to read each figure. Read Parts A--E in order the first
time. After that, use Part G (the graph guide) and Part H (the glossary)
as references. Part I has self-check questions with answers.

\begin{longtable}[]{@{}
  >{\raggedright\arraybackslash}p{(\columnwidth - 2\tabcolsep) * \real{0.5000}}
  >{\raggedright\arraybackslash}p{(\columnwidth - 2\tabcolsep) * \real{0.5000}}@{}}
\toprule\noalign{}
\begin{minipage}[b]{\linewidth}\raggedright
Part
\end{minipage} & \begin{minipage}[b]{\linewidth}\raggedright
Topic
\end{minipage} \\
\midrule\noalign{}
\endhead
\bottomrule\noalign{}
\endlastfoot
\protect\hyperlink{part-a--the-task-and-the-robots}{A} & The task and
the robots: kinematics, formations, the swept corridor, traffic \\
\protect\hyperlink{part-b--sensing-and-steering}{B} & Sensing and
steering: lidar, noise, Kalman range filter, VFH family, APF, FGM, DWA,
MPC, MPPI \\
\protect\hyperlink{part-c--safety}{C} & Safety: the collision gate, the
separating-axis test and its defect, monitors, margins, watchdog \\
\protect\hyperlink{part-d--learning-the-neural-networks}{D} & Learning:
CNNs, circular convolution, DeepSets, imitation/DAgger, oracle labels,
conformal prediction \\
\protect\hyperlink{part-e--measuring-fairly-statistics-and-protocol}{E}
& Measuring fairly: seeds, common random numbers, intervals, paired
tests, tuning, held-out evaluation \\
\protect\hyperlink{part-f--physics-and-ros-2}{F} & Physics (MuJoCo) and
ROS 2 (nodes, topics, DDS, lockstep, deployment) \\
\protect\hyperlink{part-g--how-to-read-every-graph}{G} & How to read
every graph (G1--G27) \\
\protect\hyperlink{part-h--glossary}{H} & Glossary \\
\protect\hyperlink{part-i--self-check}{I} & Self-check questions \\
\end{longtable}

Notation: \emph{r} = robot radius (0.18 m), \emph{dt} = control period
(0.1 s), σ0 = sensor range noise (standard deviation, in metres),
\emph{K} = number of robots.

\begin{center}\rule{0.5\linewidth}{0.5pt}\end{center}

\hypertarget{part-a-the-task-and-the-robots}{%
\section{Part A --- The task and the
robots}\label{part-a-the-task-and-the-robots}}

\hypertarget{a1.-the-mission}{%
\subsection{A1. The mission}\label{a1.-the-mission}}

Several small robots must carry \textbf{one rigid load} (a pallet or a
plate) from a start A to a goal B through a warehouse. Shuttles move
along fixed routes and \textbf{never give way}. A mission has three
phases:

\begin{enumerate}
\def\labelenumi{\arabic{enumi}.}
\tightlist
\item
  \textbf{Forming.} The robots start scattered. Each one drives to its
  \emph{slot}, its place under the load.
\item
  \textbf{Docking.} The robots close the last few centimetres so the
  load can be pinned on.
\item
  \textbf{Transport.} The robots move as one rigid body until the load
  reaches

  \begin{enumerate}
  \def\labelenumii{\Alph{enumii}.}
  \setcounter{enumii}{1}
  \tightlist
  \item
  \end{enumerate}
\end{enumerate}

\textbf{Success} means the load reaches B. \textbf{Safety} means no
robot or load touches a static obstacle (a rack or wall) or a mover.

\hypertarget{a2.-unicycle-kinematics-how-a-robot-moves}{%
\subsection{A2. Unicycle kinematics: how a robot
moves}\label{a2.-unicycle-kinematics-how-a-robot-moves}}

A differential-drive robot has two driven wheels. If the wheels track
their commands well, the robot behaves like a \textbf{unicycle}. It has
a position (x, y) and a heading θ, and it can only move along its
heading and turn at a limited rate:

\begin{verbatim}
θ ← θ + clip(Δθ, ±ω_max)            turn, rate-limited
x ← x + v·cos θ·dt,  y ← y + v·sin θ·dt
\end{verbatim}

It cannot slide sideways. This is called a \textbf{non-holonomic}
constraint. The control commands are a forward speed \emph{v} and a turn
rate \emph{ω}. In ROS these are the two fields of \texttt{cmd\_vel}.

\hypertarget{a3.-formations-and-the-rigid-body}{%
\subsection{A3. Formations and the rigid
body}\label{a3.-formations-and-the-rigid-body}}

\textbf{Slots.} In transport the robots are slots of one rigid frame
with centre \textbf{c} and heading θ. Robot \emph{i} sits at c +
R(θ)·offsetᵢ, where R(θ) is the 2-D rotation matrix and the offsets
define the shape: - square (4 robots); - triangle (3); - 2×3 grid (6); -
3×3 grid (9); - line.

The \textbf{spacing \emph{s}} is the distance between neighbouring
slots.

\textbf{One virtual leader.} A single leader steers the frame. Every
robot then follows its own slot. This is \textbf{leader--follower
formation control}.

\textbf{The payload body} used for collision checking is the
\textbf{convex hull} of the slots, grown by \emph{r}. The convex hull is
the smallest convex shape that contains all the slots: imagine a rubber
band around them.

\textbf{Slot assignment.} At the start, the slots are matched to robots
by the \textbf{Hungarian algorithm}, which finds the assignment with the
smallest total travel distance.

\hypertarget{a4.-the-swept-corridor-why-a-2-m-wide-body-needs-a-wider-aisle}{%
\subsection{A4. The swept corridor (why a 2 m-wide body needs a wider
aisle)}\label{a4.-the-swept-corridor-why-a-2-m-wide-body-needs-a-wider-aisle}}

The body rotates as it steers. The width it can need at some moment is
the diameter of its \textbf{circumscribed circle}, not its width when
aligned with the aisle. Some examples: - \textbf{Square at s = 1.18 m:}
the body is 1.54 m wide, but it needs a 2.03 m corridor. - \textbf{3×3
grid at s = 1.0 m:} circumradius √2 + 0.18 ≈ 1.59 m, so it needs about
3.2 m. That is why it fails in the 2.45 m corridor of the study
warehouse (E16).

\hypertarget{a5.-movers-and-the-traffic-contract}{%
\subsection{A5. Movers and the traffic
contract}\label{a5.-movers-and-the-traffic-contract}}

Movers are kinematic shuttles that travel straight segments at constant
speed and pause at the ends. Their routes are stored in a JSON
\textbf{traffic contract}. Because the motion is piecewise linear, the
future position of a mover is \emph{known exactly} from the contract.
Some algorithms use that knowledge (PA-VFH-K, and the baselines'
predicted scans). A deployed robot usually does not have it, which is
why PA-VFH does without it.

\hypertarget{a6.-traffic-phase}{%
\subsection{A6. Traffic phase}\label{a6.-traffic-phase}}

The \emph{phase} of the traffic is \emph{when} each mover starts its
route. In the original study every run used the same timetable. That is
unfair, because a planner may simply be lucky with that one timetable.
\textbf{Randomising the phase} gives each run (seed) its own start
offsets, drawn from an independent random stream, so the obstacle layout
is unchanged but the encounters differ. Part G13 shows what this did to
the results.

\hypertarget{a7.-the-mission-layer-the-same-for-every-law}{%
\subsection{A7. The mission layer: the same for every
law}\label{a7.-the-mission-layer-the-same-for-every-law}}

Only the \emph{steering vector} differs between algorithms. Everything
else is shared: - \textbf{Heading blend.} The goal direction and the
steering direction are combined,
\texttt{θ\_raw\ =\ angle(w\_goal·goal\ +\ w\_vfh·steer)}, then low-pass
filtered (α = 0.35) and rate-limited. - \textbf{Safety gate.} Every
candidate pose is checked before it is committed (Part C). -
\textbf{Escape probe.} If the gate rejects the planned pose, the probe
looks for any admissible heading nearby. - \textbf{Anti-deadlock layer.}
If progress along the route stalls for 15 s, the layer commits to a
diversified heading for a while. - \textbf{Motion limits.} These are
load-carrying limits on speed, yaw acceleration and lateral
acceleration, so the load is not yanked.

\begin{center}\rule{0.5\linewidth}{0.5pt}\end{center}

\hypertarget{part-b-sensing-and-steering}{%
\section{Part B --- Sensing and
steering}\label{part-b-sensing-and-steering}}

\hypertarget{b1.-lidar-and-tof-sensors}{%
\subsection{B1. Lidar and ToF sensors}\label{b1.-lidar-and-tof-sensors}}

A 2-D \textbf{lidar} sends out beams in a full circle (here 288 beams,
1.25° apart) and returns, for each beam, the range to the first surface
it hits. A \textbf{ToF fan} is a cheaper sensor with a narrow field of
view (60°).

\textbf{The noise model.} Each range is the true range plus Gaussian
noise with standard deviation σ0 (lidar). For ToF the noise is σ0 +
k·range. Beyond the maximum range there is ``no return''.

\hypertarget{b2.-the-kalman-range-filter-per-beam}{%
\subsection{B2. The Kalman range filter (per
beam)}\label{b2.-the-kalman-range-filter-per-beam}}

A single noisy reading is unreliable, so each beam is filtered over
time. Each beam has an \textbf{estimate} \emph{x} and a
\textbf{variance} \emph{P}. At every new reading \emph{z} with noise
variance \emph{s²}:

\begin{verbatim}
predict:   P ← P + q²                    (the world may have moved a little)
gain:      K = P / (P + s²)
update:    x ← x + K·(z − x),   P ← (1 − K)·P
reset:     if |z − x| > 3·√(P + s²), set x ← z (something new entered the beam)
\end{verbatim}

How to read it: - \textbf{K near 1:} trust the new reading. This happens
with a precise sensor. - \textbf{K near 0:} trust the history. The
filter averages a noisy sensor. - \textbf{The reset} stops the filter
lagging when a mover suddenly appears. - \textbf{√P} is the filter's own
uncertainty. It is reused downstream as σ per beam.

This one ingredient carries the histogram family's robustness to noise
(G9).

\hypertarget{b3.-the-polar-histogram-vfh-and-vfh}{%
\subsection{B3. The polar histogram: VFH and
VFH+}\label{b3.-the-polar-histogram-vfh-and-vfh}}

\textbf{VFH} (Vector Field Histogram) turns a scan into
\textbf{sectors}: 72 sectors of 5° around the robot. It then asks, for
each sector, ``is it dangerous to go this way?'':

\begin{enumerate}
\def\labelenumi{\arabic{enumi}.}
\tightlist
\item
  \textbf{Nearest return per sector:} d = the smallest range in that
  sector.
\item
  \textbf{Danger} u = exp(−d² / 2σ\_h²). Near obstacles give u ≈ 1, far
  ones u ≈ 0. The width σ\_h is set so that u crosses the
  \textbf{threshold} (0.3) at the body radius plus a reaction gap.
\item
  \textbf{VFH+ enlargement.} An obstacle at distance d blocks not only
  its own sector but every sector within the angle γ = asin(r\_body /
  d). That accounts for the size of the body.
\item
  \textbf{Smoothing} over neighbouring sectors, which wraps around at
  ±180°.
\item
  \textbf{Valleys.} Sectors below the threshold are free. Among them,
  the law picks the one closest to the goal direction, plus a
  \textbf{hysteresis} penalty for changing its mind.
\item
  \textbf{APF reflex.} A short-range push away from very close returns
  is added (B5).
\end{enumerate}

\textbf{The ladder in this project} adds one ingredient at a time: 1.
\textbf{VFH:} the basic histogram. 2. \textbf{pVFH:} adds the movers'
predicted positions over a horizon. This is the contract's future. 3.
\textbf{pVFH+F:} adds the per-beam range filter (B2). 4.
\textbf{PA-VFH-K:} adds the \emph{expected kernel} (B4). This was the
original proposal, ``P-VFH''. 5. \textbf{PA-VFH (no NN):} the same
pipeline on the robot's \emph{own instantaneous scans}, with no
contract. 6. \textbf{PA-VFH:} replaces the kernel with a learned
predictor (Part D).

\hypertarget{b4.-the-expected-kernel-and-why-it-hardly-matters}{%
\subsection{B4. The expected kernel, and why it hardly
matters}\label{b4.-the-expected-kernel-and-why-it-hardly-matters}}

If the measured range has uncertainty s, the \emph{expected} danger is

\begin{verbatim}
E[exp(−d²/2σ_h²)]  with d ~ N(d_meas, s²)  =  √(σ_h² / (σ_h² + s²)) · exp(−d_meas² / 2(σ_h² + s²))
\end{verbatim}

Noise widens the kernel, so the robot reacts a little earlier. G3 shows
that, at realistic noise, the reaction distance moves by only about 3
\%. That is why the ``uncertainty-aware'' kernel adds nothing
measurable. It is a good example of checking a mechanism's \emph{size}
before claiming it.

\hypertarget{b5.-apf-artificial-potential-field}{%
\subsection{B5. APF: artificial potential
field}\label{b5.-apf-artificial-potential-field}}

The goal attracts and obstacles repel. The repulsion from each close
return falls off as 1/range\^{}p.~APF is simple and reacts fast, but it
has well-known \textbf{local minima}: places where attraction and
repulsion cancel and the robot stops.

\hypertarget{b6.-fgm-follow-the-gap}{%
\subsection{B6. FGM: follow the gap}\label{b6.-fgm-follow-the-gap}}

FGM looks for the widest \textbf{gap} in the scan and steers towards its
middle, weighted towards the goal. It is very cheap. It can oscillate
when gaps open and close.

\hypertarget{b7.-dwa-dynamic-window-approach}{%
\subsection{B7. DWA: dynamic window
approach}\label{b7.-dwa-dynamic-window-approach}}

DWA samples (v, ω) pairs the robot can actually reach in the next
instant: the \emph{dynamic window}. It simulates each pair forward for T
seconds and scores it on heading to the goal, clearance and speed. It
picks the best pair. Because it reasons about short trajectories rather
than directions, it handles a moving body well.

\hypertarget{b8.-mpc-model-predictive-control-here-a-lattice-mpc}{%
\subsection{B8. MPC: model predictive control (here, a lattice
MPC)}\label{b8.-mpc-model-predictive-control-here-a-lattice-mpc}}

MPC predicts the future with a model and optimises a cost over a
horizon. It applies the first step, then re-plans every step
(\textbf{receding horizon}). The lattice MPC here: - builds a small tree
of heading changes over two segments; - scores each branch on progress,
a smooth obstacle cost (soft safety distance \texttt{d\_safe}) and a
hard clearance; - keeps the best branch.

\texttt{d\_safe} keeps MPC away from walls, which is why it survived
physics without a margin (G21).

\hypertarget{b9.-mppi-model-predictive-path-integral}{%
\subsection{B9. MPPI: model predictive path
integral}\label{b9.-mppi-model-predictive-path-integral}}

MPPI is a sampling-based MPC: 1. Sample N = 256 noisy control sequences
around the previous plan. 2. Roll each one out and compute its cost. 3.
Average them with weights exp(−cost/λ).

λ (the temperature) sets how greedy the averaging is. MPPI handles
non-convex costs well, at higher compute (about 4 ms per step here).

\hypertarget{b10.-the-information-each-law-uses}{%
\subsection{B10. The information each law
uses}\label{b10.-the-information-each-law-uses}}

\begin{longtable}[]{@{}
  >{\raggedright\arraybackslash}p{(\columnwidth - 4\tabcolsep) * \real{0.3333}}
  >{\raggedright\arraybackslash}p{(\columnwidth - 4\tabcolsep) * \real{0.3333}}
  >{\raggedright\arraybackslash}p{(\columnwidth - 4\tabcolsep) * \real{0.3333}}@{}}
\toprule\noalign{}
\begin{minipage}[b]{\linewidth}\raggedright
Law
\end{minipage} & \begin{minipage}[b]{\linewidth}\raggedright
Reads the movers' future (contract)?
\end{minipage} & \begin{minipage}[b]{\linewidth}\raggedright
Notes
\end{minipage} \\
\midrule\noalign{}
\endhead
\bottomrule\noalign{}
\endlastfoot
VFH, pVFH, pVFH+F, PA-VFH-K, FGM, DWA, MPC, MPPI & yes, through the
predicted scan & the benchmark's standard sensing front end \\
PA-VFH (no NN), PA-VFH & \textbf{no}: own instantaneous scans only &
deployable \\
NN (DAgger) & scan features as for MPC & learned \\
\end{longtable}

So PA-VFH competes with \emph{less information} than the baselines. Keep
this in mind when it loses on contacts with movers.

\begin{center}\rule{0.5\linewidth}{0.5pt}\end{center}

\hypertarget{part-c-safety}{%
\section{Part C --- Safety}\label{part-c-safety}}

\hypertarget{c1.-the-safety-gate-atomic-pose-commit}{%
\subsection{C1. The safety gate (atomic pose
commit)}\label{c1.-the-safety-gate-atomic-pose-commit}}

Before the frame moves, the gate checks the \textbf{next pose}: the body
at the new centre and heading, at t and t + dt. If any part overlaps an
obstacle, the move is \textbf{rejected} and the escape logic runs.
Nothing is committed half-way, which is why the commit is called atomic.

In \texttt{truth} mode the gate uses the simulator's exact geometry.
That is a \textbf{screening device}: it measures what a law asks for
while an oracle prevents static collisions. It is not something a real
robot has.

\hypertarget{c2.-the-separating-axis-test-sat-and-the-defect-that-was-found}{%
\subsection{C2. The separating-axis test (SAT), and the defect that was
found}\label{c2.-the-separating-axis-test-sat-and-the-defect-that-was-found}}

\textbf{The theorem.} Two convex shapes do not overlap if and only if
there is a line (axis) onto which their projections do not overlap. For
boxes it is enough to test the axes perpendicular to each box's edges.

\textbf{The defect.} The MATLAB reference listed a box's corners in the
order (+,+), (+,−), (−,+), (−,−). That is \textbf{not cyclic}. Joining
consecutive corners gives two sides and two diagonals, so one of the
box's two edge normals was never tested. The test then sometimes found
``no separating axis'' for a pose that was clear: a \textbf{false
rejection}.

\textbf{Consequences.} The defect is \textbf{conservative}: it never
accepted a real collision, it only refused some safe poses. It inflated
rejection counts and lengthened some missions. Fixing the order to
cyclic, (+,+), (+,−), (−,−), (−,+), made MATLAB and Python agree in all
221 grid cells (G7).

\hypertarget{c3.-ground-truth-screening-vs-deployable-monitors}{%
\subsection{C3. Ground-truth screening vs deployable
monitors}\label{c3.-ground-truth-screening-vs-deployable-monitors}}

\begin{longtable}[]{@{}
  >{\raggedright\arraybackslash}p{(\columnwidth - 6\tabcolsep) * \real{0.2500}}
  >{\raggedright\arraybackslash}p{(\columnwidth - 6\tabcolsep) * \real{0.2500}}
  >{\raggedright\arraybackslash}p{(\columnwidth - 6\tabcolsep) * \real{0.2500}}
  >{\raggedright\arraybackslash}p{(\columnwidth - 6\tabcolsep) * \real{0.2500}}@{}}
\toprule\noalign{}
\begin{minipage}[b]{\linewidth}\raggedright
Gate source
\end{minipage} & \begin{minipage}[b]{\linewidth}\raggedright
Who decides each pose
\end{minipage} & \begin{minipage}[b]{\linewidth}\raggedright
What it needs
\end{minipage} & \begin{minipage}[b]{\linewidth}\raggedright
Used for
\end{minipage} \\
\midrule\noalign{}
\endhead
\bottomrule\noalign{}
\endlastfoot
\texttt{truth} & exact geometry & the simulator & reference behaviour \\
\texttt{audit} & truth decides; the sensor monitor is scored in its
shadow & the simulator & measuring the monitor without changing the
mission \\
\texttt{sensor} & the geometric monitor & only the robots' scans &
deployable \\
\texttt{learned} & the learned monitor with a conformal margin & only
the scans & deployable \\
\end{longtable}

\textbf{The geometric monitor} converts the filtered scan into points.
It admits a pose only if every point stays clear of the body by

\begin{verbatim}
r + k·σ  (sensor uncertainty, k = 3)  +  v_obs·h  (a point may move up to v_obs = 1 m/s for h = 0.2 s)
\end{verbatim}

\textbf{Two error types:} - \textbf{False accept:} the monitor admitted
a pose that the truth rejects. This is dangerous. - \textbf{False
reject:} the monitor refused a safe pose. This costs time.

\hypertarget{c4.-physical-safety-margin-gate_margin}{%
\subsection{\texorpdfstring{C4. Physical safety margin
(\texttt{gate\_margin})}{C4. Physical safety margin (gate\_margin)}}\label{c4.-physical-safety-margin-gate_margin}}

A kinematic gate with r = 0.18 m lets a planner pass 1 mm from a rack. A
real or physically simulated robot lags its command by centimetres, so
it touches the rack. \texttt{gate\_margin\ =\ 0.10\ m} makes the gate
demand r + 0.10 m. It is applied \textbf{to every law} in batches P3 and
P4. G21 shows it working.

\hypertarget{c5.-the-ros-2-watchdog}{%
\subsection{C5. The ROS 2 watchdog}\label{c5.-the-ros-2-watchdog}}

A software watchdog in the controller node does three things: - If any
robot's odometry or scan is older than 0.5 s, it \textbf{stops every
robot} and publishes \texttt{/swarm/safety\ =\ 1}. - It clamps every
command to \texttt{max\_v} and \texttt{max\_w}. - It sends zero velocity
on shutdown.

Measured over DDS, the team stopped 0.6 s after a lidar failure (G26).
It complements a hardware e-stop; it is not a replacement for one.

\begin{center}\rule{0.5\linewidth}{0.5pt}\end{center}

\hypertarget{part-d-learning-the-neural-networks}{%
\section{Part D --- Learning: the neural
networks}\label{part-d-learning-the-neural-networks}}

\hypertarget{d1.-what-a-neural-network-is-in-this-projects-terms}{%
\subsection{D1. What a neural network is, in this project's
terms}\label{d1.-what-a-neural-network-is-in-this-projects-terms}}

A neural network is a function with many tunable numbers (weights). It
is fitted to examples by gradient descent on a \textbf{loss}. The
networks here are small, about 11 k weights for the occupancy model.
They are trained with PyTorch and \textbf{run with numpy only}, so a
robot needs no GPU and no deep-learning library.

\hypertarget{d2.-convolution-and-why-it-is-circular}{%
\subsection{\texorpdfstring{D2. Convolution, and why it is
\emph{circular}}{D2. Convolution, and why it is circular}}\label{d2.-convolution-and-why-it-is-circular}}

A \textbf{1-D convolution} slides a small filter (5 sectors wide here)
along a sequence. It detects the same local pattern wherever it appears.

A scan is a \textbf{ring}: sector 71 neighbours sector 0. A
\emph{circular} convolution wraps around, so the network treats every
bearing the same way. Rotating the scan by k sectors rotates the output
by k sectors. This property is called \textbf{rotation equivariance}. It
saves data and generalises better.

\hypertarget{d3.-the-pa-vfh-occupancy-network-input-output-training}{%
\subsection{D3. The PA-VFH occupancy network: input, output,
training}\label{d3.-the-pa-vfh-occupancy-network-input-output-training}}

\textbf{Input:} a 6 × 72 array (6 channels per sector): - three channels
for the last three filtered scans, each holding the sector's minimum
range minus the body radius; - the range uncertainty σ per sector; - an
observed/unobserved mask; - the formation's own closing speed towards
each sector.

\textbf{Network:} 3 circular convolution layers (32 channels, width 5,
ReLU), then a 1×1 output layer and a sigmoid.

\textbf{Output:} a danger value u ∈ (0, 1) per sector, \emph{for the
next 0.5 s}.

\textbf{Labels (oracle):} the simulator knows where every mover will be,
so it can compute the \emph{true} danger over the next 0.5 s exactly.
That is the training target. Training uses privileged information that
the robot never has at run time, sometimes called ``learning by
cheating''. The network learns to infer from three scans what the oracle
knows.

\textbf{Loss:} binary cross-entropy with soft labels (the targets are
numbers between 0 and 1, not just 0 or 1). Sectors that really are
dangerous get more weight.

\textbf{On-policy rounds:} 1. Round 1 collects data while the
non-learned PA-VFH drives. 2. Round 2 collects data while PA-VFH itself
drives, and retrains on the union.

This fixes \textbf{distribution shift}: the learned planner visits
states the first driver never did.

\textbf{Where it sits:} it replaces only the danger kernel. Valley
selection and the gate remain classical, which keeps the system
auditable.

\hypertarget{d4.-deepsets-and-the-team-size-agnostic-dagger-policy}{%
\subsection{D4. DeepSets, and the team-size-agnostic DAgger
policy}\label{d4.-deepsets-and-the-team-size-agnostic-dagger-policy}}

The \textbf{DAgger policy} steers the whole formation directly. - It
encodes the scan with a circular CNN. - It encodes the slots with
\textbf{DeepSets}: apply the same small network to each slot, then
\emph{sum}. A sum does not care about order or count, so one policy
works for 3, 4, 6 or 9 robots.

\hypertarget{d5.-imitation-learning-and-dagger}{%
\subsection{D5. Imitation learning and
DAgger}\label{d5.-imitation-learning-and-dagger}}

\textbf{Behaviour cloning} copies an expert (here MPC) from recorded
data. Its weakness is that once the learner makes a small mistake, it
reaches states the expert never showed it.

\textbf{DAgger} (Dataset Aggregation) fixes this: 1. Let the
\emph{learner} drive. 2. Ask the \emph{expert} what it would have done
in the states the learner visited. 3. Add those labels to the dataset
and retrain. 4. Repeat.

\hypertarget{d6.-conformal-prediction-a-calibrated-error-margin}{%
\subsection{D6. Conformal prediction: a calibrated error
margin}\label{d6.-conformal-prediction-a-calibrated-error-margin}}

The learned monitor regresses the clearance. A regression can be wrong,
and split conformal prediction turns it into a decision with a stated
error rate: 1. On a \textbf{calibration set} not used for training,
compute the residuals \emph{e = true − predicted}. 2. Take \emph{q\_α} =
the ⌈(n+1)(1−α)⌉-th smallest residual. For α = 0.05, that is roughly the
95th percentile. 3. \textbf{Admit} a pose only if \emph{predicted − q\_α
\textgreater{} r}.

\textbf{Guarantee:} if new checks are drawn like the calibration ones
(exchangeable), the true clearance falls below the corrected estimate at
most α of the time. The guarantee is \textbf{marginal} (on average over
checks) and \textbf{in-distribution only}. On fused per-robot scans,
which it was never calibrated on, it admitted more unsafe poses
(ROS2\_VALIDATION.md §2).

\hypertarget{d7.-two-ways-to-judge-a-learned-predictor}{%
\subsection{D7. Two ways to judge a learned
predictor}\label{d7.-two-ways-to-judge-a-learned-predictor}}

\begin{itemize}
\tightlist
\item
  \textbf{Value error} (mean absolute error against the oracle): how
  close the numbers are.
\item
  \textbf{Decision agreement:} do the network and the oracle agree on
  \emph{blocked or free} at the 0.3 threshold?
\end{itemize}

Steering depends on the second. In G23 the network has a \emph{larger}
value error (0.022 vs 0.012) but \emph{better} decisions (98.2 \% vs
98.0 \%). It puts small, hedged values on free sectors that do not
change any decision.

\begin{center}\rule{0.5\linewidth}{0.5pt}\end{center}

\hypertarget{part-e-measuring-fairly-statistics-and-protocol}{%
\section{Part E --- Measuring fairly: statistics and
protocol}\label{part-e-measuring-fairly-statistics-and-protocol}}

\hypertarget{e1.-monte-carlo-runs-seeds-and-cells}{%
\subsection{E1. Monte Carlo runs, seeds and
cells}\label{e1.-monte-carlo-runs-seeds-and-cells}}

One mission is one random outcome, so each configuration (a
\textbf{cell}: world × law × settings) is run with several
\textbf{seeds}, usually 6--10. Results are reported per cell or pooled.

\hypertarget{e2.-common-random-numbers-crn}{%
\subsection{E2. Common random numbers
(CRN)}\label{e2.-common-random-numbers-crn}}

All laws in a cell use the \emph{same} seeds, so they face the same
noise samples and the same traffic. Differences between laws are then
not caused by luck. This is what makes \textbf{paired} comparisons
possible.

\hypertarget{e3.-rates-and-the-wilson-interval}{%
\subsection{E3. Rates and the Wilson
interval}\label{e3.-rates-and-the-wilson-interval}}

A table entry like \texttt{69/80\ {[}0.77,\ 0.92{]}} means 69 successes
out of 80 runs, with a \textbf{95 \% Wilson interval} of 0.77--0.92 for
the true rate. The Wilson interval is preferred over ``p ±
1.96·√(p(1−p)/n)'' because it behaves well near 0 and 1 and for small
\emph{n}. With 10 runs and 10/10 successes, the interval is {[}0.72,
1.00{]}: a perfect score on 10 runs is weak evidence.

\hypertarget{e4.-means-and-the-bootstrap-interval}{%
\subsection{E4. Means and the bootstrap
interval}\label{e4.-means-and-the-bootstrap-interval}}

For averages (time, contact steps), the \textbf{bootstrap} resamples the
runs with replacement 2,000 times and recomputes the mean each time. The
middle 95 \% of those means is the interval. It needs no normality
assumption.

\hypertarget{e5.-paired-differences-and-not-significant}{%
\subsection{E5. Paired differences and ``not
significant''}\label{e5.-paired-differences-and-not-significant}}

In E19 each mission (world, noise, seed) is run \emph{with} and
\emph{without} the network. The difference per mission is computed, and
then its bootstrap interval, for example −19.3 {[}−57.9, +0.7{]}
rejected steps. - \textbf{Negative} favours the network. - \textbf{The
interval touches 0}, so the data are consistent with no effect. The
honest wording is ``consistent but not statistically significant''.

\hypertarget{e6.-equal-budget-tuning-and-held-out-evaluation}{%
\subsection{E6. Equal-budget tuning and held-out
evaluation}\label{e6.-equal-budget-tuning-and-held-out-evaluation}}

\textbf{Equal-budget tuning:} every law gets the same search, 24
configurations including its defaults, on the same 7 training missions.
This stops ``tuned proposal vs untuned baseline'' comparisons.

\textbf{Held-out evaluation:} the comparison then runs on worlds and
seeds the tuner never saw (E18). Training and tuning used seeds ≥ 101;
evaluation used seeds 1--10.

\hypertarget{e7.-cross-stack-agreement}{%
\subsection{E7. Cross-stack agreement}\label{e7.-cross-stack-agreement}}

The same experiment was run in MATLAB and in Python. A cell
\textbf{agrees} if the 95 \% intervals of reach, rejections and minimum
gap overlap. 221/221 cells agreed. At zero noise the runs are compared
step by step, and the largest centroid difference was 7 × 10⁻⁷ m.

\hypertarget{e8.-the-metrics}{%
\subsection{E8. The metrics}\label{e8.-the-metrics}}

\begin{longtable}[]{@{}
  >{\raggedright\arraybackslash}p{(\columnwidth - 2\tabcolsep) * \real{0.5000}}
  >{\raggedright\arraybackslash}p{(\columnwidth - 2\tabcolsep) * \real{0.5000}}@{}}
\toprule\noalign{}
\begin{minipage}[b]{\linewidth}\raggedright
Metric
\end{minipage} & \begin{minipage}[b]{\linewidth}\raggedright
Meaning
\end{minipage} \\
\midrule\noalign{}
\endhead
\bottomrule\noalign{}
\endlastfoot
reach & the load got to B within the time limit \\
contact steps / contact runs & steps, or runs, with any overlap with a
static or a mover (ground truth) \\
min gap & the smallest clearance minus \emph{r}. Negative means
overlap \\
rejections & steps the gate refused \\
time & mission time of the runs that reached B \\
comp ms & the steering law's compute per step \\
clean run (physics) & reached, with no static and no mover contact in
MuJoCo \\
false accept / reject & monitor errors, scored against the truth \\
\end{longtable}

\begin{center}\rule{0.5\linewidth}{0.5pt}\end{center}

\hypertarget{part-f-physics-and-ros-2}{%
\section{Part F --- Physics and ROS 2}\label{part-f-physics-and-ros-2}}

\hypertarget{f1.-why-physics}{%
\subsection{F1. Why physics}\label{f1.-why-physics}}

The kinematic model assumes perfect tracking. \textbf{MuJoCo} simulates
masses, wheel friction, motor torque limits, contacts and the load
pinned to each robot by a ball joint. The robots now \textbf{lag} and
\textbf{slip}: - the \emph{formation error} (robot-to-slot distance)
grows in turns; - planners that drive close to obstacles touch them.

There are three \textbf{payload modes}: - \textbf{virtual:} no physical
load; - \textbf{on top:} a plate and crate pinned to the robots; -
\textbf{behind:} a towed sled. The towed load swings \emph{outside} the
formation footprint the planner models, so every law struggles with it.

Signals measured in physics: - \textbf{pin gap:} how far the joint
between robot and load is pulled apart; - \textbf{tilt:} the load's
tilt; - \textbf{contact force.}

\hypertarget{f2.-the-robot-parameter-sweep-model-mismatch}{%
\subsection{F2. The robot-parameter sweep (model
mismatch)}\label{f2.-the-robot-parameter-sweep-model-mismatch}}

P2 and P4 change one physical parameter at a time while the planner
keeps its nominal model: - chassis mass ×0.5 / ×2; - wheel speed 0.8 /
1.6 m/s; - friction 0.6; - payload ×0.5 / ×2; - motor force ×0.5.

This measures robustness to a robot that is not what the planner thinks.

\hypertarget{f3.-ros-2-in-one-page}{%
\subsection{F3. ROS 2 in one page}\label{f3.-ros-2-in-one-page}}

\begin{itemize}
\tightlist
\item
  \textbf{Node:} a process with a role. Here \texttt{mujoco\_plant}
  (physics), \texttt{formation\_controller} (mission and steering),
  \texttt{mover\_node} and \texttt{swarm\_logger}.
\item
  \textbf{Topic:} a named message stream, such as
  \texttt{/robot1/cmd\_vel} (geometry\_msgs/Twist),
  \texttt{/robot1/odom} (nav\_msgs/Odometry), \texttt{/robot1/scan}
  (sensor\_msgs/LaserScan), and \texttt{/swarm/phase} and
  \texttt{/swarm/safety} (std\_msgs/Int32).
\item
  \textbf{DDS:} the middleware that carries messages between processes
  and machines. \textbf{QoS} settings choose reliable or best-effort
  delivery and queue depth.
\item
  \textbf{Launch file:} starts several nodes with parameters, for
  example
  \texttt{ros2\ launch\ ros2\_swarm\_transport\ robots.launch.py}.
\item
  \textbf{colcon:} the build tool that installs the package
  (executables, launch files, worlds, model weights).
\end{itemize}

\hypertarget{f4.-lockstep-vs-free-running}{%
\subsection{F4. Lockstep vs
free-running}\label{f4.-lockstep-vs-free-running}}

\begin{itemize}
\tightlist
\item
  \textbf{Lockstep (validation):}

  \begin{enumerate}
  \def\labelenumi{\arabic{enumi}.}
  \tightlist
  \item
    The plant publishes the state.
  \item
    The controller replies with commands.
  \item
    Only then does physics advance one tick.
  \end{enumerate}

  The run is deterministic, so it can be compared to the in-process run
  \emph{exactly} (G25).
\item
  \textbf{Free-running (deployment):} the controller steps on a 10 Hz
  timer using the latest \texttt{/odom} and \texttt{/scan}. That is what
  real robots do (G26).
\end{itemize}

\hypertarget{f4b.-simulation-time-use_sim_time}{%
\subsection{\texorpdfstring{F4b. Simulation time
(\texttt{use\_sim\_time})}{F4b. Simulation time (use\_sim\_time)}}\label{f4b.-simulation-time-use_sim_time}}

A free-running controller steps on a timer. Two kinds of clock exist: -
\textbf{The wall clock.} If the simulator falls behind real time (for
example while rendering video), the controller runs ahead of the
physics. Its watchdog then sees ``stale'' data, and the two drift apart.
- \textbf{Simulation time.} With \texttt{use\_sim\_time:=true}, the
timer and the watchdog follow \texttt{/clock}, which the MuJoCo plant
publishes, so both sides stay in step whatever the simulator's speed.

This project measured the difference. On the wall clock, a rendered run
timed out with contacts; on sim time, the same run was clean
(ROS2\_VALIDATION.md §1c).

\hypertarget{f4c.-physics-through-ros-2-r1-and-r2}{%
\subsection{F4c. Physics through ROS 2: R1 and
R2}\label{f4c.-physics-through-ros-2-r1-and-r2}}

\begin{itemize}
\tightlist
\item
  \textbf{R1:} the P3 physics batch with the plant and controller as ROS
  2 processes in lockstep. All 90 runs are identical to the in-process
  runs in the same environment.
\item
  \textbf{Host vs container:} different MuJoCo versions differ in about
  the 14th digit. In a chaotic mission, one flipped gate decision then
  changes the whole outcome. That is why identity is checked within one
  environment, and why environments are compared statistically.
\item
  \textbf{R2:} the deployment stack, with the plant free-running and
  publishing \texttt{/odom} and \texttt{/scan} from the physical robots.
\end{itemize}

\hypertarget{f5.-real-sensor-mode-from-laserscans-to-a-formation-scan}{%
\subsection{F5. Real-sensor mode: from LaserScans to a formation
scan}\label{f5.-real-sensor-mode-from-laserscans-to-a-formation-scan}}

In deployment the controller must not read the simulator's map of
obstacles. Instead: 1. Each robot's \texttt{LaserScan} is converted to
points in the map frame, using that robot's odometry pose. 2.
\textbf{Self-filtering:} points on the team-mates, and under the carried
load, are removed. The filter radius grows with noise, r + 0.10 + 3σ. 3.
A \textbf{virtual scan} from the formation centre is built from the
fused points. Each beam takes the \emph{median of its three nearest
points}, not the minimum. The minimum of \emph{n} noisy points is biased
short by about σ·√(2 ln n). 4. PA-VFH and the monitor then work exactly
as in simulation.

Only PA-VFH (with \texttt{gate\_source} sensor or learned) is allowed in
this mode. The node refuses configurations that need the oracle.

\begin{center}\rule{0.5\linewidth}{0.5pt}\end{center}

\hypertarget{part-g-how-to-read-every-graph}{%
\section{Part G --- How to read every
graph}\label{part-g-how-to-read-every-graph}}

For each figure: \textbf{what is plotted}, \textbf{how to read it},
\textbf{what it shows in this project}, and \textbf{pitfalls}.

\textbf{General rules:} - \textbf{Error bars} are 95 \% intervals:
Wilson for rates, bootstrap for means. If two intervals don't overlap,
the difference is very likely real. If they overlap, look for a
\emph{paired} comparison before concluding anything. - \textbf{Count the
runs.} ``10/10'' on 10 runs is weaker evidence than ``130/150''. -
\textbf{Check which gate was used.} In \texttt{truth} mode, static
contacts mean the \emph{physical} or \emph{mover} side failed, not the
oracle.

\hypertarget{g1.-formation-shapes-docsfiguresconcept_formations.png}{%
\subsection{\texorpdfstring{G1. Formation shapes
(\texttt{docs/figures/concept\_formations.png})}{G1. Formation shapes (docs/figures/concept\_formations.png)}}\label{g1.-formation-shapes-docsfiguresconcept_formations.png}}

\includegraphics[width=0.95\textwidth,height=\textheight]{/home/user/Cooperative-payload-transport-by-a-four-robot-formation/docs/figures/concept_formations.png}

\begin{itemize}
\tightlist
\item
  \textbf{Plotted:} the four shapes. Disks are robots, the shaded area
  is the collision body (hull + r), the dashed circle is the
  circumscribed circle.
\item
  \textbf{Read:} the corridor a shape needs is the \emph{circle's
  diameter} (A4). Compare it with the aisle widths in G12.
\item
  \textbf{Shows:} why the line formation needs more room than its width
  suggests, and why the large 3×3 payload fails in a 2.45 m corridor.
\end{itemize}

\hypertarget{g2.-the-polar-histogram-concept_histogram.png}{%
\subsection{\texorpdfstring{G2. The polar histogram
(\texttt{concept\_histogram.png})}{G2. The polar histogram (concept\_histogram.png)}}\label{g2.-the-polar-histogram-concept_histogram.png}}

\includegraphics[width=0.95\textwidth,height=\textheight]{/home/user/Cooperative-payload-transport-by-a-four-robot-formation/docs/figures/concept_histogram.png}

\begin{itemize}
\tightlist
\item
  \textbf{Plotted:}

  \begin{itemize}
  \tightlist
  \item
    \textbf{Left:} one noisy scan (blue beams) from the formation
    centre, a mover (orange) and the body outline.
  \item
    \textbf{Right:} the 72-sector histogram. Grey bars are above the
    threshold (blocked), green bars are below it (free). The dashed
    green line is the goal bearing; the purple line is the chosen
    sector.
  \end{itemize}
\item
  \textbf{Read:} if the goal sector is green, the law goes nearly
  straight. Otherwise it picks the nearest green valley.
\item
  \textbf{Pitfall:} the enlargement (B3) blocks \emph{more} sectors than
  the obstacle itself covers. Wide grey regions near a narrow aisle are
  why histogram laws hesitate at aisle mouths (G11).
\end{itemize}

\hypertarget{g3.-the-expected-kernel-concept_kernel.png}{%
\subsection{\texorpdfstring{G3. The expected kernel
(\texttt{concept\_kernel.png})}{G3. The expected kernel (concept\_kernel.png)}}\label{g3.-the-expected-kernel-concept_kernel.png}}

\includegraphics[width=0.95\textwidth,height=\textheight]{/home/user/Cooperative-payload-transport-by-a-four-robot-formation/docs/figures/concept_kernel.png}

\begin{itemize}
\tightlist
\item
  \textbf{Plotted:}

  \begin{itemize}
  \tightlist
  \item
    \textbf{Left:} danger against measured range for four noise levels.
  \item
    \textbf{Right:} the range where danger crosses the threshold (the
    reaction distance) against noise. The green band is the study's
    noise range.
  \end{itemize}
\item
  \textbf{Read:} inside the green band the curve barely moves, 2.01 m to
  about 2.07 m.
\item
  \textbf{Shows:} the whole effect of the ``uncertainty-aware'' kernel
  is about 3 \%. That predicts, correctly, that PA-VFH-K ≈ pVFH+F in
  experiments (G9).
\end{itemize}

\hypertarget{g4.-the-range-filter-concept_range_filter.png}{%
\subsection{\texorpdfstring{G4. The range filter
(\texttt{concept\_range\_filter.png})}{G4. The range filter (concept\_range\_filter.png)}}\label{g4.-the-range-filter-concept_range_filter.png}}

\includegraphics[width=0.95\textwidth,height=\textheight]{/home/user/Cooperative-payload-transport-by-a-four-robot-formation/docs/figures/concept_range_filter.png}

\begin{itemize}
\tightlist
\item
  \textbf{Plotted:} one beam over time. Black is the true range, orange
  dots are measurements, purple is the filter estimate with ±2σ.
\item
  \textbf{Read:}

  \begin{itemize}
  \tightlist
  \item
    Left (precise sensor): the purple line follows the dots, gain ≈ 1.
  \item
    Right (noisy sensor): the purple line is smooth.
  \item
    At 4.5 s a mover enters, and the reset makes the estimate jump at
    once instead of lagging.
  \end{itemize}
\item
  \textbf{Shows:} why pVFH+F is so much more robust than pVFH (G9).
\end{itemize}

\hypertarget{g5.-the-gate-defect-concept_gate.png}{%
\subsection{\texorpdfstring{G5. The gate defect
(\texttt{concept\_gate.png})}{G5. The gate defect (concept\_gate.png)}}\label{g5.-the-gate-defect-concept_gate.png}}

\includegraphics[width=0.95\textwidth,height=\textheight]{/home/user/Cooperative-payload-transport-by-a-four-robot-formation/docs/figures/concept_gate.png}

\begin{itemize}
\tightlist
\item
  \textbf{Plotted:} the same pose under two tests.

  \begin{itemize}
  \tightlist
  \item
    \textbf{Left:} the reference ``edges'' in red, two sides plus two
    diagonals. The pose is reported as overlapping although it is 0.29 m
    clear.
  \item
    \textbf{Right:} with cyclic corners, all four sides are tested
    (green) and the pose is accepted.
  \end{itemize}
\item
  \textbf{Read:} a missing axis can only cause false \emph{rejections}
  (C2).
\end{itemize}

\hypertarget{g6.-zero-noise-cross-check-xcheck_trajectories_fixedgate.png}{%
\subsection{\texorpdfstring{G6. Zero-noise cross-check
(\texttt{xcheck\_trajectories\_fixedgate.png})}{G6. Zero-noise cross-check (xcheck\_trajectories\_fixedgate.png)}}\label{g6.-zero-noise-cross-check-xcheck_trajectories_fixedgate.png}}

\includegraphics[width=0.95\textwidth,height=\textheight]{/home/user/Cooperative-payload-transport-by-a-four-robot-formation/docs/figures/xcheck_trajectories_fixedgate.png}

\begin{itemize}
\tightlist
\item
  \textbf{Plotted:}

  \begin{itemize}
  \tightlist
  \item
    \textbf{Top:} the centroid path of each law. The wide pale line is
    MATLAB, the thin line Python. Grey is static, red dashed is the
    mover routes.
  \item
    \textbf{Bottom:} the per-step position difference.
  \end{itemize}
\item
  \textbf{Read:} if the thin line sits exactly on the wide one and the
  difference stays near zero, the two implementations are the same
  algorithm.
\item
  \textbf{Shows:} 7 × 10⁻⁷ m maximum, for every law. That is the
  rounding of the MATLAB logs.
\item
  \textbf{Pitfall:} this holds only at zero noise. With noise the random
  numbers differ, so you compare statistics instead (G7).
\end{itemize}

\hypertarget{g7.-cross-stack-scatter-cross_stack_scatter_fixedgate.png}{%
\subsection{\texorpdfstring{G7. Cross-stack scatter
(\texttt{cross\_stack\_scatter\_fixedgate.png})}{G7. Cross-stack scatter (cross\_stack\_scatter\_fixedgate.png)}}\label{g7.-cross-stack-scatter-cross_stack_scatter_fixedgate.png}}

\includegraphics[width=0.95\textwidth,height=\textheight]{/home/user/Cooperative-payload-transport-by-a-four-robot-formation/docs/figures/cross_stack_scatter_fixedgate.png}

\begin{itemize}
\tightlist
\item
  \textbf{Plotted:} each dot is one experiment cell. x is the MATLAB
  mean, y the Python mean, the line is y = x. There are three panels:
  time, rejections (log scale) and minimum gap.
\item
  \textbf{Read:}

  \begin{itemize}
  \tightlist
  \item
    Dots on the line mean the two stacks measured the same thing.
  \item
    Blue means the intervals overlap.
  \item
    A red diamond means a disagreement.
  \end{itemize}
\item
  \textbf{Shows:} with the corrected gate, 221/221 cells agree. The
  correlations are 0.98 (time), 0.95 (rejections) and 0.99 (gap).
\item
  \textbf{Pitfall:} the spread at high rejection counts is expected.
  Those cells are chaotic failures, where a single seed changes the mean
  a lot.
\end{itemize}

\hypertarget{g8.-e1-both-stacks-e1_both_stacks_fixedgate.png}{%
\subsection{\texorpdfstring{G8. E1 both stacks
(\texttt{e1\_both\_stacks\_fixedgate.png})}{G8. E1 both stacks (e1\_both\_stacks\_fixedgate.png)}}\label{g8.-e1-both-stacks-e1_both_stacks_fixedgate.png}}

\includegraphics[width=0.95\textwidth,height=\textheight]{/home/user/Cooperative-payload-transport-by-a-four-robot-formation/docs/figures/e1_both_stacks_fixedgate.png}

\begin{itemize}
\tightlist
\item
  \textbf{Plotted:} each law's metrics with 95 \% whiskers. Blue is
  MATLAB, orange Python.
\item
  \textbf{Read:} paired bars at the same height mean agreement.
\item
  \textbf{Pitfall:} APF and FGM have erratic missions. One timeout among
  10 runs shifts their means, so the whiskers are long.
\end{itemize}

\hypertarget{g9.-the-ablation-ladder-against-noise-e2_noise_both_stacks_fixedgate.png}{%
\subsection{\texorpdfstring{G9. The ablation ladder against noise
(\texttt{e2\_noise\_both\_stacks\_fixedgate.png})}{G9. The ablation ladder against noise (e2\_noise\_both\_stacks\_fixedgate.png)}}\label{g9.-the-ablation-ladder-against-noise-e2_noise_both_stacks_fixedgate.png}}

\includegraphics[width=0.95\textwidth,height=\textheight]{/home/user/Cooperative-payload-transport-by-a-four-robot-formation/docs/figures/e2_noise_both_stacks_fixedgate.png}

\begin{itemize}
\tightlist
\item
  \textbf{Plotted:} one column per law (VFH, pVFH, pVFH+F, PA-VFH-K).
  The top row is time, the bottom row is rejections, both against sensor
  noise σ0.
\item
  \textbf{Read left to right:} each column adds one ingredient. The
  ingredient that matters is the one where the curves flatten.
\item
  \textbf{Shows:}

  \begin{itemize}
  \tightlist
  \item
    VFH → pVFH: no change in robustness.
  \item
    pVFH → pVFH+F: much flatter. At σ0 = 0.1, time 77 → 53 s and
    rejections about 40 → 3.
  \item
    pVFH+F → PA-VFH-K: no difference.
  \end{itemize}
\item
  \textbf{Lesson:} the filter carries the robustness. The kernel adds
  nothing.
\end{itemize}

\hypertarget{g10.-gate-effect-gate_effect.png}{%
\subsection{\texorpdfstring{G10. Gate effect
(\texttt{gate\_effect.png})}{G10. Gate effect (gate\_effect.png)}}\label{g10.-gate-effect-gate_effect.png}}

\includegraphics[width=0.95\textwidth,height=\textheight]{/home/user/Cooperative-payload-transport-by-a-four-robot-formation/docs/figures/gate_effect.png}

\begin{itemize}
\tightlist
\item
  \textbf{Plotted:} the same seeds with the reference gate (grey) and
  the exact gate (purple), pooled over several experiments.
\item
  \textbf{Read:} the pooled means move little, because cells where laws
  genuinely fail dominate them.
\item
  \textbf{Shows:} 71--91 \% of runs are bit-identical under the two
  gates. The defect shows only where a falsely refused pose occurred.
\end{itemize}

\hypertarget{g11.-warehouse-paths-warehouse_paths.png-warehouse_matlab_vs_python.png}{%
\subsection{\texorpdfstring{G11. Warehouse paths
(\texttt{warehouse\_paths.png},
\texttt{warehouse\_matlab\_vs\_python.png})}{G11. Warehouse paths (warehouse\_paths.png, warehouse\_matlab\_vs\_python.png)}}\label{g11.-warehouse-paths-warehouse_paths.png-warehouse_matlab_vs_python.png}}

\includegraphics[width=0.95\textwidth,height=\textheight]{/home/user/Cooperative-payload-transport-by-a-four-robot-formation/docs/figures/warehouse_paths.png}

\includegraphics[width=0.95\textwidth,height=\textheight]{/home/user/Cooperative-payload-transport-by-a-four-robot-formation/docs/figures/warehouse_matlab_vs_python.png}

\begin{itemize}
\tightlist
\item
  \textbf{Plotted:} centroid paths of each law from the dock to the
  aisle. Left is without motion limits, right with them.
\item
  \textbf{Read:}

  \begin{itemize}
  \tightlist
  \item
    Smooth paths mean decisive steering.
  \item
    Zig-zags at aisle mouths mean histogram hesitation (G2 pitfall).
  \end{itemize}
\item
  \textbf{The second figure} overlays MATLAB and Python on the
  validation warehouse, a world neither implementation was developed
  against. They agree again to 7 × 10⁻⁷ m.
\end{itemize}

\hypertarget{g12.-generated-layouts-resultsfigures_extfig_layouts.png}{%
\subsection{\texorpdfstring{G12. Generated layouts
(\texttt{results/figures\_ext/fig\_layouts.png})}{G12. Generated layouts (results/figures\_ext/fig\_layouts.png)}}\label{g12.-generated-layouts-resultsfigures_extfig_layouts.png}}

\includegraphics[width=0.95\textwidth,height=\textheight]{/home/user/Cooperative-payload-transport-by-a-four-robot-formation/results/figures_ext/fig_layouts.png}

\begin{itemize}
\tightlist
\item
  \textbf{Plotted:} the six generated warehouses (30×20 to 80×50 m).
  Racks and columns are dark grey. The blue line is the global A* route,
  from the start (blue dot) to the goal (green star). The dashed orange
  lines are the mover lanes. Each title gives the size and the number of
  statics and movers.
\item
  \textbf{Read:} compare the aisle width with the swept diameter of the
  formation (G1). \texttt{wh\_m\_narrow} has 2.4 m aisles.
\end{itemize}

\hypertarget{g13.-traffic-phase-effect-fig_phase_effect.png}{%
\subsection{\texorpdfstring{G13. Traffic-phase effect
(\texttt{fig\_phase\_effect.png})}{G13. Traffic-phase effect (fig\_phase\_effect.png)}}\label{g13.-traffic-phase-effect-fig_phase_effect.png}}

\includegraphics[width=0.95\textwidth,height=\textheight]{/home/user/Cooperative-payload-transport-by-a-four-robot-formation/results/figures_ext/fig_phase_effect.png}

\begin{itemize}
\tightlist
\item
  \textbf{Plotted:} for each law, two markers per panel. The blue circle
  is the nominal timetable, the orange square is random phase, each with
  a 95 \% Wilson bar.

  \begin{itemize}
  \tightlist
  \item
    \textbf{Left:} fraction of runs reaching the goal.
  \item
    \textbf{Right:} fraction of runs with any contact.
  \end{itemize}
\item
  \textbf{Read:} horizontal distance between circle and square is the
  effect of randomising traffic. Bars that don't overlap mean a real
  effect.
\item
  \textbf{Shows:}

  \begin{itemize}
  \tightlist
  \item
    For every law, contact runs rise sharply. MPC goes from 19 \% to 47
    \%.
  \item
    Reach drops 5--15 points.
  \item
    The fixed timetable had flattered every planner, MPC most of all,
    because its clean result depended on timing.
  \end{itemize}
\item
  \textbf{Pitfall:} the ranking among laws also changes. Results from a
  fixed timetable should not be used to rank planners.
\end{itemize}

\hypertarget{g14.-e15-reach-per-layout-fig_e15_reach.png}{%
\subsection{\texorpdfstring{G14. E15 reach per layout
(\texttt{fig\_e15\_reach.png})}{G14. E15 reach per layout (fig\_e15\_reach.png)}}\label{g14.-e15-reach-per-layout-fig_e15_reach.png}}

\includegraphics[width=0.95\textwidth,height=\textheight]{/home/user/Cooperative-payload-transport-by-a-four-robot-formation/results/figures_ext/fig_e15_reach.png}

\begin{itemize}
\tightlist
\item
  \textbf{Plotted:} a heatmap. Rows are layouts, columns are laws, each
  cell is ``reached / runs''. Darker means a higher success rate.
\item
  \textbf{Read:}

  \begin{itemize}
  \tightlist
  \item
    Look for \textbf{light rows}: layouts where every law fails.
  \item
    Look for \textbf{light columns}: laws that fail everywhere.
  \end{itemize}
\item
  \textbf{Shows:}

  \begin{itemize}
  \tightlist
  \item
    The \texttt{wh\_m\_narrow} row is 0/10 for every law. That is an
    open problem, not a weakness of one law.
  \item
    The study warehouse (``dynamic'') is the hardest of the passable
    layouts.
  \item
    PA-VFH and NN reach 69/80, the best.
  \end{itemize}
\end{itemize}

\hypertarget{g15.-e16-team-size-fig_e16_dynamic.png-fig_e16_wh_m_racks.png-fig_e16_wh_l_racks.png}{%
\subsection{\texorpdfstring{G15. E16 team size
(\texttt{fig\_e16\_dynamic.png}, \texttt{fig\_e16\_wh\_m\_racks.png},
\texttt{fig\_e16\_wh\_l\_racks.png})}{G15. E16 team size (fig\_e16\_dynamic.png, fig\_e16\_wh\_m\_racks.png, fig\_e16\_wh\_l\_racks.png)}}\label{g15.-e16-team-size-fig_e16_dynamic.png-fig_e16_wh_m_racks.png-fig_e16_wh_l_racks.png}}

\includegraphics[width=0.95\textwidth,height=\textheight]{/home/user/Cooperative-payload-transport-by-a-four-robot-formation/results/figures_ext/fig_e16_dynamic.png}

\includegraphics[width=0.95\textwidth,height=\textheight]{/home/user/Cooperative-payload-transport-by-a-four-robot-formation/results/figures_ext/fig_e16_wh_m_racks.png}

\includegraphics[width=0.95\textwidth,height=\textheight]{/home/user/Cooperative-payload-transport-by-a-four-robot-formation/results/figures_ext/fig_e16_wh_l_racks.png}

\begin{itemize}
\tightlist
\item
  \textbf{Plotted:} a heatmap per world. Rows are teams (shape, K,
  spacing), columns are laws, each cell is ``reached / 10''.
\item
  \textbf{Read:} compare rows in the \emph{same} world.
\item
  \textbf{Shows:}

  \begin{itemize}
  \tightlist
  \item
    The 3×3 grid at s = 1.0 (a 2 × 2 m payload, about 3.2 m swept
    diameter) fails in the study warehouse for every law except MPPI
    (7/10). It passes in the large rack warehouses.
  \item
    K = 9 at s = 0.59 (the same small payload with more robots) is easy.
  \item
    So it is \textbf{payload size relative to the corridor}, not robot
    count, that limits.
  \item
    NN (DAgger) at 2×3, K = 6 in the study warehouse (2/10) shows a
    learned policy's weak spot on a team it was trained on but in a hard
    world.
  \end{itemize}
\end{itemize}

\hypertarget{g16.-e17-safety-monitors-fig_e17_monitor.png}{%
\subsection{\texorpdfstring{G16. E17 safety monitors
(\texttt{fig\_e17\_monitor.png})}{G16. E17 safety monitors (fig\_e17\_monitor.png)}}\label{g16.-e17-safety-monitors-fig_e17_monitor.png}}

\includegraphics[width=0.95\textwidth,height=\textheight]{/home/user/Cooperative-payload-transport-by-a-four-robot-formation/results/figures_ext/fig_e17_monitor.png}

\begin{itemize}
\tightlist
\item
  \textbf{Plotted:} three gate types (ground truth, sensor monitor,
  learned + conformal) × two noise levels (circle 0.02, square 0.10).

  \begin{itemize}
  \tightlist
  \item
    \textbf{Left:} reach rate.
  \item
    \textbf{Right:} contact-run rate.
  \end{itemize}
\item
  \textbf{Read:}

  \begin{itemize}
  \tightlist
  \item
    Compare each monitor with ground truth at the same noise.
  \item
    A monitor that keeps reach but adds contacts is less safe.
  \item
    A monitor that loses reach is too conservative.
  \end{itemize}
\item
  \textbf{Shows:}

  \begin{itemize}
  \tightlist
  \item
    The sensor monitor matches ground truth at σ0 = 0.02 but collapses
    at 0.10 (reach about 0.5). Its 3σ inflation closes the aisles.
  \item
    The learned monitor keeps reach at both noise levels.
  \item
    Both deployable monitors have more contact runs than ground truth.
    The contacts are mostly movers driving into the formation, which no
    pose check can prevent.
  \end{itemize}
\end{itemize}

\hypertarget{g17.-robot-platform-and-mission-frames-robot_platform.png-mujoco_forming_sequence.png-mujoco_frames.png}{%
\subsection{\texorpdfstring{G17. Robot platform and mission frames
(\texttt{robot\_platform.png}, \texttt{mujoco\_forming\_sequence.png},
\texttt{mujoco\_frames.png})}{G17. Robot platform and mission frames (robot\_platform.png, mujoco\_forming\_sequence.png, mujoco\_frames.png)}}\label{g17.-robot-platform-and-mission-frames-robot_platform.png-mujoco_forming_sequence.png-mujoco_frames.png}}

\includegraphics[width=0.95\textwidth,height=\textheight]{/home/user/Cooperative-payload-transport-by-a-four-robot-formation/docs/figures/robot_platform.png}

\includegraphics[width=0.95\textwidth,height=\textheight]{/home/user/Cooperative-payload-transport-by-a-four-robot-formation/docs/figures/mujoco_forming_sequence.png}

\includegraphics[width=0.95\textwidth,height=\textheight]{/home/user/Cooperative-payload-transport-by-a-four-robot-formation/docs/figures/mujoco_frames.png}

\begin{itemize}
\tightlist
\item
  \textbf{Plotted:} the MuJoCo robot (radius 0.18 m, 6.9 kg, two wheels
  plus casters), the square formation with and without the load, and
  frames from the mission: scattered start, forming, docking, transport.
\item
  \textbf{Read:} use these to picture what the kinematic numbers mean
  physically.
\end{itemize}

\hypertarget{g18.-physics-tracking-mujoco_tracking.png}{%
\subsection{\texorpdfstring{G18. Physics tracking
(\texttt{mujoco\_tracking.png})}{G18. Physics tracking (mujoco\_tracking.png)}}\label{g18.-physics-tracking-mujoco_tracking.png}}

\includegraphics[width=0.95\textwidth,height=\textheight]{/home/user/Cooperative-payload-transport-by-a-four-robot-formation/docs/figures/mujoco_tracking.png}

\begin{itemize}
\tightlist
\item
  \textbf{Plotted:}

  \begin{itemize}
  \tightlist
  \item
    \textbf{Left:} the planner's kinematic frame (dashed), the physical
    robots (green) and the load (orange).
  \item
    \textbf{Top right:} the worst slot error over time per payload mode.
  \item
    \textbf{Bottom right:} the frame's yaw rate.
  \end{itemize}
\item
  \textbf{Read:} error peaks line up with yaw-rate peaks. Turns are
  where physics departs from the kinematic model.
\end{itemize}

\hypertarget{g19.-physics-summary-mujoco_summary.png}{%
\subsection{\texorpdfstring{G19. Physics summary
(\texttt{mujoco\_summary.png})}{G19. Physics summary (mujoco\_summary.png)}}\label{g19.-physics-summary-mujoco_summary.png}}

\includegraphics[width=0.95\textwidth,height=\textheight]{/home/user/Cooperative-payload-transport-by-a-four-robot-formation/docs/figures/mujoco_summary.png}

\begin{itemize}
\tightlist
\item
  \textbf{Plotted:} formation error and contact steps by law and payload
  mode.
\item
  \textbf{Read:} the towed (``behind'') mode has the largest errors for
  every law (F1).
\end{itemize}

\hypertarget{g20.-motion-limits-mujoco_motion_limits.png}{%
\subsection{\texorpdfstring{G20. Motion limits
(\texttt{mujoco\_motion\_limits.png})}{G20. Motion limits (mujoco\_motion\_limits.png)}}\label{g20.-motion-limits-mujoco_motion_limits.png}}

\includegraphics[width=0.95\textwidth,height=\textheight]{/home/user/Cooperative-payload-transport-by-a-four-robot-formation/docs/figures/mujoco_motion_limits.png}

\begin{itemize}
\tightlist
\item
  \textbf{Plotted:} the worst slot error over time with the
  load-carrying motion limits off (red) and on (green). Left MPC, right
  PA-VFH-K.
\item
  \textbf{Read:} with limits on, the error peaks are lower. The limits
  stop the planner asking for turns the robots cannot follow with a
  load.
\end{itemize}

\hypertarget{g21.-s1-the-physical-margin-resultsfigures_signalss1_physics_margin_clearance.png}{%
\subsection{\texorpdfstring{G21. S1: the physical margin
(\texttt{results/figures\_signals/S1\_physics\_margin\_clearance.png})}{G21. S1: the physical margin (results/figures\_signals/S1\_physics\_margin\_clearance.png)}}\label{g21.-s1-the-physical-margin-resultsfigures_signalss1_physics_margin_clearance.png}}

\includegraphics[width=0.95\textwidth,height=\textheight]{/home/user/Cooperative-payload-transport-by-a-four-robot-formation/results/figures_signals/S1_physics_margin_clearance.png}

\begin{itemize}
\tightlist
\item
  \textbf{Plotted:}

  \begin{itemize}
  \tightlist
  \item
    \textbf{Top:} the planner's clearance to static obstacles during
    transport in MuJoCo, same seed and traffic. Red is PA-VFH without a
    margin, blue PA-VFH with the 0.10 m margin, amber MPC. The dashed
    red line is contact.
  \item
    \textbf{Bottom:} the time steps where physics recorded a rack
    contact.
  \end{itemize}
\item
  \textbf{Read:} contacts in the bottom strip coincide with the red
  curve touching zero (around 33--46 s).
\item
  \textbf{Shows:}

  \begin{itemize}
  \tightlist
  \item
    Without a margin, PA-VFH hugs the racks.
  \item
    With the margin it keeps ≥ 0.1 m and the contacts disappear.
  \item
    MPC keeps a wide berth by itself (its \texttt{d\_safe}).
  \end{itemize}
\item
  \textbf{Batch result:} PA-VFH clean 0/20 → 16/20 (carried and virtual
  load).
\end{itemize}

\hypertarget{g22.-s2-one-physics-mission-signal-by-signal-s2_physics_signals.png}{%
\subsection{\texorpdfstring{G22. S2: one physics mission, signal by
signal
(\texttt{S2\_physics\_signals.png})}{G22. S2: one physics mission, signal by signal (S2\_physics\_signals.png)}}\label{g22.-s2-one-physics-mission-signal-by-signal-s2_physics_signals.png}}

\includegraphics[width=0.95\textwidth,height=\textheight]{/home/user/Cooperative-payload-transport-by-a-four-robot-formation/results/figures_signals/S2_physics_signals.png}

\begin{itemize}
\tightlist
\item
  \textbf{Plotted:} PA-VFH with the margin, load on top. The panels from
  top:

  \begin{itemize}
  \tightlist
  \item
    payload speed;
  \item
    formation error;
  \item
    payload tilt;
  \item
    pin gap;
  \item
    clearance;
  \item
    an event strip with the transport phase, mover contact, static
    contact and gate rejections.
  \end{itemize}
\item
  \textbf{Read across panels at the same time:}

  \begin{itemize}
  \tightlist
  \item
    When speed drops (a slow-down near traffic), the formation error
    rises and gate rejections cluster.
  \item
    The tilt stays below 0.04° and the pin gap below 1 mm: the load is
    level and attached.
  \end{itemize}
\item
  \textbf{Shows:} no contacts at all, and a stable load throughout.
\end{itemize}

\hypertarget{g23.-s3-inside-the-neural-network-s3_pavfh_danger_maps.png}{%
\subsection{\texorpdfstring{G23. S3: inside the neural network
(\texttt{S3\_pavfh\_danger\_maps.png})}{G23. S3: inside the neural network (S3\_pavfh\_danger\_maps.png)}}\label{g23.-s3-inside-the-neural-network-s3_pavfh_danger_maps.png}}

\includegraphics[width=0.95\textwidth,height=\textheight]{/home/user/Cooperative-payload-transport-by-a-four-robot-formation/results/figures_signals/S3_pavfh_danger_maps.png}

\begin{itemize}
\tightlist
\item
  \textbf{Plotted:} three heatmaps of danger. x is time, y is bearing
  (−180° to 180°), colour is danger (white = free, dark blue = blocked).

  \begin{enumerate}
  \def\labelenumi{\arabic{enumi}.}
  \tightlist
  \item
    The oracle (true danger over the next 0.5 s).
  \item
    The PA-VFH network, which sees only three noisy scans.
  \item
    The analytic kernel, from the current scan only.
  \end{enumerate}

  The bottom panel shows each method's mean error against the oracle
  over time.
\item
  \textbf{Read:}

  \begin{itemize}
  \tightlist
  \item
    Diagonal stripes are obstacles sweeping past as the formation moves
    or turns.
  \item
    Compare the shapes in the network's map with the oracle's. A good
    predictor shows the same stripes, slightly ahead of where the
    current scan has them.
  \end{itemize}
\item
  \textbf{Shows:}

  \begin{itemize}
  \tightlist
  \item
    The network reproduces the oracle's structure closely.
  \item
    Decision agreement is 98.2 \% (network) vs 98.0 \% (kernel); missed
    blocked sectors are 8.0 \% vs 8.3 \%.
  \item
    Its value error is higher, 0.022 vs 0.012, because of hedged small
    values (D7).
  \item
    There is a \textbf{start-up transient} at 5--10 s, before three
    frames of history exist, where it over-predicts danger everywhere.
  \end{itemize}
\item
  \textbf{Pitfall:} don't judge the network by value error alone.
  Steering uses only blocked vs free.
\end{itemize}

\hypertarget{g24.-s4-the-sensor-monitor-against-ground-truth-s4_monitor_signals.png}{%
\subsection{\texorpdfstring{G24. S4: the sensor monitor against ground
truth
(\texttt{S4\_monitor\_signals.png})}{G24. S4: the sensor monitor against ground truth (S4\_monitor\_signals.png)}}\label{g24.-s4-the-sensor-monitor-against-ground-truth-s4_monitor_signals.png}}

\includegraphics[width=0.95\textwidth,height=\textheight]{/home/user/Cooperative-payload-transport-by-a-four-robot-formation/results/figures_signals/S4_monitor_signals.png}

\begin{itemize}
\tightlist
\item
  \textbf{Plotted:}

  \begin{itemize}
  \tightlist
  \item
    \textbf{Left:} every gate check. x is the true clearance minus r, y
    is the monitor's clearance minus its required gap. Blue points were
    admitted, orange rejected. The shaded area left of x = 0 is truly
    unsafe.
  \item
    \textbf{Right:} both signals over the sequence of checks.
  \end{itemize}
\item
  \textbf{Read:}

  \begin{itemize}
  \tightlist
  \item
    A \textbf{false accept} is a blue point in the shaded area.
  \item
    Points below the diagonal mean the monitor is \emph{more
    pessimistic} than the truth.
  \item
    Orange points with x \textgreater{} 0 are false rejects, which cost
    time.
  \end{itemize}
\item
  \textbf{Shows:} the monitor stays below the truth (conservative) and
  made 0 false accepts in this mission.
\end{itemize}

\hypertarget{g25.-s5-ros-2-lockstep-equality-s5_ros2_lockstep_equality.png}{%
\subsection{\texorpdfstring{G25. S5: ROS 2 lockstep equality
(\texttt{S5\_ros2\_lockstep\_equality.png})}{G25. S5: ROS 2 lockstep equality (S5\_ros2\_lockstep\_equality.png)}}\label{g25.-s5-ros-2-lockstep-equality-s5_ros2_lockstep_equality.png}}

\includegraphics[width=0.95\textwidth,height=\textheight]{/home/user/Cooperative-payload-transport-by-a-four-robot-formation/results/figures_signals/S5_ros2_lockstep_equality.png}

\begin{itemize}
\tightlist
\item
  \textbf{Plotted:}

  \begin{itemize}
  \tightlist
  \item
    \textbf{Left:} the centroid path from two ROS 2 processes over DDS
    (thin blue) on top of the in-process path (wide amber).
  \item
    \textbf{Right:} their difference over time.
  \end{itemize}
\item
  \textbf{Read:} a flat line at 0 means identical.
\item
  \textbf{Shows:} a 0 m difference over 1,010 steps. The ROS 2 wiring
  (messages, packing, timing protocol) adds no error.
\end{itemize}

\hypertarget{g26.-s6-real-sensor-deployment-over-dds-s6_ros2_deployment_signals.png}{%
\subsection{\texorpdfstring{G26. S6: real-sensor deployment over DDS
(\texttt{S6\_ros2\_deployment\_signals.png})}{G26. S6: real-sensor deployment over DDS (S6\_ros2\_deployment\_signals.png)}}\label{g26.-s6-real-sensor-deployment-over-dds-s6_ros2_deployment_signals.png}}

\includegraphics[width=0.95\textwidth,height=\textheight]{/home/user/Cooperative-payload-transport-by-a-four-robot-formation/results/figures_signals/S6_ros2_deployment_signals.png}

\begin{itemize}
\tightlist
\item
  \textbf{Plotted:} four panels.

  \begin{enumerate}
  \def\labelenumi{\arabic{enumi}.}
  \tightlist
  \item
    A histogram of the time between \texttt{cmd\_vel} messages.
  \item
    robot1's commanded speed.
  \item
    The true gap to the nearest obstacle, which the controller never
    sees.
  \item
    The watchdog test: commanded speed (blue) and \texttt{/swarm/safety}
    (red) around a lidar failure at 20 s.
  \end{enumerate}
\item
  \textbf{Read:}

  \begin{itemize}
  \tightlist
  \item
    Panel 1 should be a sharp peak at 100 ms.
  \item
    Panel 3 dipping below zero means contact. Here the dips are movers
    driving into the formation; there are no static contacts.
  \item
    Panel 4: the flag rises and the speed drops to zero shortly after
    the failure.
  \end{itemize}
\item
  \textbf{Shows:}

  \begin{itemize}
  \tightlist
  \item
    Cadence: median 100 ms, 95th percentile 104 ms.
  \item
    3/3 missions reached the goal from scans alone.
  \item
    The watchdog stopped the team 0.6 s after the failure.
  \end{itemize}
\end{itemize}

\hypertarget{g27.-tables-and-the-demos-bar-charts}{%
\subsection{G27. Tables and the demo's bar
charts}\label{g27.-tables-and-the-demos-bar-charts}}

\begin{itemize}
\tightlist
\item
  \texttt{k/n\ {[}lo,\ hi{]}} is a rate with a Wilson interval (E3).
  \texttt{m\ {[}lo,\ hi{]}} is a mean with a bootstrap interval (E4).
\item
  In R9, the ``paired'' rows are per-mission differences (E5).
\item
  The demo's bar charts show raw counts. Read them together with the run
  counts in the axis label.
\end{itemize}

\begin{center}\rule{0.5\linewidth}{0.5pt}\end{center}

\hypertarget{part-h-glossary}{%
\section{Part H --- Glossary}\label{part-h-glossary}}

\begin{longtable}[]{@{}
  >{\raggedright\arraybackslash}p{(\columnwidth - 2\tabcolsep) * \real{0.5000}}
  >{\raggedright\arraybackslash}p{(\columnwidth - 2\tabcolsep) * \real{0.5000}}@{}}
\toprule\noalign{}
\begin{minipage}[b]{\linewidth}\raggedright
Term
\end{minipage} & \begin{minipage}[b]{\linewidth}\raggedright
Meaning
\end{minipage} \\
\midrule\noalign{}
\endhead
\bottomrule\noalign{}
\endlastfoot
ADL & anti-deadlock layer (A7) \\
APF & artificial potential field (B5) \\
audit mode & truth decides; the monitor is scored in its shadow (C3) \\
bootstrap & resampling interval for a mean (E4) \\
circular convolution & convolution that wraps around a ring of sectors
(D2) \\
clean run & reached, with no contact in physics (E8) \\
conformal prediction & calibrated error margin for a learned predictor
(D6) \\
contract & JSON file of mover routes and timing (A5) \\
CRN & common random numbers (E2) \\
DAgger & dataset aggregation for imitation learning (D5) \\
DDS & ROS 2 middleware (F3) \\
DeepSets & permutation-invariant set encoder: per-element net, then sum
(D4) \\
DWA & dynamic window approach (B7) \\
enlargement (VFH+) & blocking neighbouring sectors by the body size
(B3) \\
false accept / reject & monitor admitted an unsafe / refused a safe pose
(C3) \\
FGM & follow-the-gap method (B6) \\
gate & pose check before every move (C1) \\
gate\_margin & extra clearance the gate demands (C4) \\
held-out & worlds and seeds not used in training or tuning (E6) \\
hysteresis & penalty for switching direction (B3) \\
lockstep & plant and controller advance tick by tick (F4) \\
MPC / MPPI & model predictive control / its sampling-based variant (B8,
B9) \\
oracle labels & training targets computed from privileged simulator
knowledge (D3) \\
PA-VFH & the proposed law: VFH+ with a learned occupancy predictor \\
PA-VFH (no NN) & the same pipeline with the analytic kernel (the
ablation) \\
PA-VFH-K & former ``P-VFH'': analytic kernel plus contract prediction \\
phase (traffic) & when each mover starts (A6) \\
pin gap & separation at the robot--load joint (F1) \\
SAT & separating-axis theorem (C2) \\
self-filter & removing team and load returns from fused scans (F5) \\
sector & 5° bin of the polar histogram (B3) \\
σ0 & sensor range noise, standard deviation (B1) \\
slot & a robot's place in the formation (A3) \\
swept corridor & width the rotating body can need (A4) \\
unicycle & kinematic model of a differential-drive robot (A2) \\
VFH / VFH+ & vector field histogram / with enlargement (B3) \\
watchdog & stops all robots on stale sensor data (C5) \\
Wilson interval & interval for a rate (E3) \\
\end{longtable}

\begin{center}\rule{0.5\linewidth}{0.5pt}\end{center}

\hypertarget{part-i-self-check}{%
\section{Part I --- Self-check}\label{part-i-self-check}}

Answers follow each question.

\begin{enumerate}
\def\labelenumi{\arabic{enumi}.}
\item
  \textbf{Why does the square formation at s = 1.18 m need a 2.03 m
  corridor when it is only 1.54 m wide?} Because it rotates. The widest
  extent over all headings is the circumscribed circle's diameter (A4).
\item
  \textbf{The reference gate never caused a collision. Why was fixing it
  still important?} It falsely rejected safe poses. That inflated
  rejections, lengthened missions, and broke agreement between the two
  stacks in one cell. A benchmark must measure the planners, not a test
  artefact (C2, G7).
\item
  \textbf{Which ingredient of the histogram ladder makes it robust to
  noise, and which figure shows it?} The per-beam Kalman range filter
  (pVFH → pVFH+F). Figure E2 (G9).
\item
  \textbf{Why does randomising the traffic phase matter more than adding
  seeds with the same timetable?} With a fixed timetable every seed
  meets the same encounters, so seeds only vary sensor noise. Planners
  can look safe through lucky timing. G13 shows contact rates roughly
  doubling once the timing varies.
\item
  \textbf{E19 reports −19.3 {[}−57.9, +0.7{]} rejected steps for the
  network. What can you claim?} The network tends to reduce rejections
  (negative is good), but the interval touches 0. The evidence is
  consistent with no effect, so the claim is ``consistent but not
  significant'' (E5). The pre-registered E20, with 360 fresh pairs,
  confirmed this: the rejection effect vanished (+0.12 {[}−2.71,
  +3.00{]}). Only mission time improved significantly (−0.97 s {[}−1.66,
  −0.33{]}). That is why a small exploratory result needs a confirmatory
  run before it is claimed.
\item
  \textbf{In S3 the network has a larger mean error than the analytic
  kernel but is called better. Contradiction?} No.~Steering depends on
  blocked vs free decisions, where the network agrees more often with
  the oracle. Its larger value error comes from small values on free
  sectors that never cross the threshold (D7).
\item
  \textbf{What is the difference between ``oracle labels'' and ``the
  contract''?} Both use knowledge of the movers' future. PA-VFH-K uses
  the contract \emph{at run time}, which is not deployable. PA-VFH uses
  future knowledge only \emph{to make training labels}. At run time it
  sees only scans (D3, B10).
\item
  \textbf{The learned monitor's conformal guarantee says α = 5 \%. Why
  did it admit more unsafe poses in ROS 2 real-sensor mode?} The
  guarantee holds for checks distributed like the calibration set. Fused
  per-robot scans are a different distribution (D6).
\item
  \textbf{Why did MPC survive physics without a margin while PA-VFH did
  not?} MPC's cost keeps a soft safety distance (\texttt{d\_safe}) from
  obstacles. The histogram laws accept any pose the gate allows, which
  can be 1 mm from a rack (G21).
\item
  \textbf{In S6, panel 3 dips below zero but there are no static
  contacts. What happened?} Movers never yield and drove into the
  formation. The controller cannot prevent a collision that is entirely
  the mover's doing (A1, G26).
\item
  \textbf{Why must the self-filter grow with sensor noise?} Returns from
  team-mates scatter by about 3σ beyond the robot's radius. A fixed
  filter leaves them in the cloud as phantom obstacles next to the
  slots, and forming never finishes (F5).
\item
  \textbf{What sensor would you specify for a deployment, and why?} A
  2-D lidar with range noise ≤ 5 cm. Above that, the 3σ inflation plus
  the robot radius nearly fills the gap between formation and wall in
  the narrowest corridor, and reach collapses (ROS2\_VALIDATION §2b).
\end{enumerate}

\end{document}
